\chapter{Introdução}\label{cap:intro}


O receio de que as máquinas tornem o trabalho humano dispensável é recorrente na história econômica, mas poucas vezes voltou ao centro do debate público com a intensidade observada após o lançamento do ChatGPT, em 30 de novembro de 2022. A difusão acelerada dos modelos generativos de linguagem reabriu uma pergunta antiga sob uma forma nova: se as ondas anteriores de automação atingiam sobretudo tarefas rotineiras e manuais, a \gls{ia} contemporânea avança sobre tarefas \emph{cognitivas}, tais como redigir, programar, analisar e sintetizar, antes tidas como reduto seguro do trabalho qualificado. Essa inquietação ganhou um lastro empírico com as estimativas de que parcela expressiva das ocupações é tecnicamente suscetível à computadorização \parencite{frey2017} e, mais recentemente, de que cerca de 80\% da força de trabalho americana está em ocupações com ao menos um décimo das tarefas expostas aos grandes modelos de linguagem, e quase um quinto em ocupações com ao menos metade das tarefas expostas, inclusive as de alta remuneração \parencite{eloundou2024}. Tal preocupação tornou-se objeto de uma literatura econômica em expansão acelerada.

Essa literatura, contudo, não é uníssona. De um lado, a tradição da economia do trabalho construída desde os anos 2000 -- a abordagem de tarefas e o arcabouço de \textcite{acemogluRestrepo2019} -- ensina que o efeito de uma tecnologia sobre o emprego não é um veredito fixo, mas o saldo entre forças opostas: o deslocamento de tarefas, a criação de novas tarefas, a complementaridade que amplia a produtividade e a expansão da demanda agregada.

De outro, a evidência empírica inicial sobre a IA generativa é heterogênea e por vezes surpreendente. Experimentos encontram ganhos de produtividade concentrados nos trabalhadores menos experientes \parencite{noyZhang2023, brynjolfssonLiRaymond2025}, sugerindo complementaridade onde se temia substituição.

Ao deslocar, porém, o nível de análise, encontramos resultados distintos. No plano agregado, \textcite{kauhanenRouvinen2025}, com dados populacionais da Finlândia, não detectaram no curto prazo diferenças significativas de emprego ou de salário entre as ocupações mais e menos expostas à IA generativa.  Há, portanto, um descompasso entre os ganhos medidos em experimentos individuais e a ausência de efeitos macroeconômicos relevantes, o que reforça, agora para os modelos generativos, o achado de \textcite{acemogluVacancies2022} sobre a IA preditiva. Ao mesmo tempo, quando se observam as respostas comportamentais, emergem efeitos antecipatórios. \textcite{gollerGschwendtWolter2025} documentaram queda acentuada na procura por formação profissional nas ocupações cognitivas mais expostas, o que reforça a evidência de que o mercado reage à ameaça percebida antes mesmo de seus efeitos se materializarem no emprego.

Além de diferenças de magnitude, encontramos também divergências em relação à direção dos efeitos. A evidência experimental localiza os maiores ganhos -- e, por implicação, a maior pressão substitutiva -- na base da distribuição de qualificações, ao passo que os índices de exposição ocupacional e os estudos de escolha de carreira apontam o risco para as tarefas cognitivas de alta qualificação \parencite{eloundou2024, gollerGschwendtWolter2025}; e há ainda quem mostre que a própria adoção da tecnologia permanece restringida por trabalho humano qualificado e escasso \parencite{fontanelli2025}. São diagnósticos que divergem, como se argumenta no Capítulo~\ref{cap:discussao}, menos por erro do que por medirem objetos distintos -- produtividade, emprego, salário ou escolha ocupacional --, em unidades e horizontes também distintos.

O resultado é um campo vasto, disperso e de diagnósticos não convergentes, em que coexistem a tese da polarização herdada da era da automação e a conjectura de que a IA generativa deslocaria o risco para o topo da distribuição de qualificações. É essa tensão que motiva a presente revisão.

A pergunta central que organiza este trabalho é: \emph{como o avanço da IA generativa alterou o diagnóstico da literatura econômica sobre o impacto da inteligência artificial no mercado de trabalho, e em que medida essa alteração acompanha o objeto tecnológico estudado, e não apenas o calendário?} A hipótese de trabalho que orienta esse exame é a de que a literatura publicada antes do lançamento do ChatGPT, que tendia a situar o risco da automação nas tarefas rotineiras de \emph{baixa} qualificação, foi sucedida por uma literatura que desloca esse risco para as tarefas cognitivas de \emph{alta} qualificação -- uma hipótese de ruptura no enquadramento do problema. Subordinam-se a ela perguntas auxiliares sobre quais mecanismos teóricos a literatura mobiliza, como mede os efeitos, em que horizonte temporal se situa e qual o sinal predominante que atribui ao saldo de empregos, e sobre se esses traços se alteram de forma sistemática entre os dois períodos.

O objetivo geral é \textbf{mapear e sintetizar sistematicamente a literatura econômica sobre os impactos da \gls{ia} no mercado de trabalho publicada a partir de 2013 e indexada até maio de 2026, comparando o enquadramento do problema antes e depois da difusão da IA generativa}. Desse objetivo decorrem os seguintes objetivos específicos:

\begin{enumerate}
  \item caracterizar o corpus de estudos quanto a volume, evolução temporal,
    tipo de evidência, tecnologia de IA em foco e atributos metodológicos;
  \item ler o corpus longitudinalmente em três janelas tecnológicas --- a da
    automação (\janelaUm{}), a do aprendizado de máquina e \textit{deep learning}
    (\janelaDois{}) e a da IA generativa (\janelaTres{}) ---, identificando como
    variam os mecanismos teóricos invocados;
  \item comparar, dimensão a dimensão, os estudos anteriores e posteriores ao
    lançamento do ChatGPT, testando a hipótese de deslocamento do risco para a
    alta qualificação;
  \item distinguir continuidades e rupturas entre os dois períodos e avaliar a
    robustez dos achados a definições alternativas de período, a subconjuntos
    do corpus e à composição de tipos de estudo;
  \item separar, usando os estudos sobre robótica e automação como grupo de
    comparação, as mudanças associadas ao objeto tecnológico das tendências
    gerais do campo.
\end{enumerate}

A justificativa do trabalho é dupla. Em primeiro lugar, há uma lacuna de síntese. A produção sobre IA e trabalho cresceu de forma acelerada e fragmentada, com desenhos empíricos heterogêneos -- exercícios de exposição ocupacional, estudos macro e setoriais, experimentos, análises de firmas e de plataformas de trabalho \textit{online} -- e vereditos que não convergem. Falta uma revisão sistemática que organize esse material sob uma lente teórica única e estabeleça, com método explícito e replicável, como e em que medida o diagnóstico do campo se transformou com a IA generativa e o que dos paradigmas da era da automação se manteve. Esta revisão se propõe a preencher essa lacuna.

Em segundo lugar, há uma justificativa metodológica. A comparação pareada pré/pós-ChatGPT, ancorada nos quatro mecanismos do arcabouço de tarefas e aplicada a um corpus de centenas de estudos com extração padronizada, oferece um protocolo transparente para distinguir mudança de ênfase de mudança das conclusões dos trabalhos, algo que a leitura narrativa convencional não permite. Esse protocolo é narrativo-estruturado, e não meta-analítico, uma vez que a heterogeneidade de desenhos e a baixa comparabilidade das magnitudes reportadas inviabilizam a agregação estatística de efeitos.

Os resultados, antecipados aqui em linhas gerais, respondem à pergunta central nos seguintes termos. Como a codificação registra o grupo de qualificação \emph{atingido} pela tecnologia, considera-se que um estudo situa um grupo em risco apenas quando não atribui sinal positivo ao efeito sobre o emprego. Assim medida, a fração de estudos que situam o risco nas ocupações de alta qualificação multiplica-se por quatro entre os períodos, de 2,0\% para 8,5\% dos que classificam o grupo atingido, ainda que o risco para a baixa qualificação permaneça o diagnóstico mais frequente. Essa mudança está associada sobretudo ao \emph{objeto} tecnológico: entre os estudos sobre IA generativa, o risco para a alta qualificação é a categoria mais frequente (15 de 33 estudos, 45,5\%), ao passo que, entre os estudos sobre robótica e automação, ele é residual nos dois períodos. Controlado o objeto, o efeito do período de publicação deixa de se distinguir de zero, embora o pequeno número de estudos não permita descartar um efeito residual de calendário. A literatura não mudou de ideia sobre a automação; passou a estudar também uma tecnologia com outro perfil de tarefas, padrão compatível com a abordagem de tarefas. Uma ressalva acompanha o achado: parte das demais mudanças observadas -- o horizonte mais curto, o recuo dos mecanismos de demanda agregada, o sinal mais frequentemente positivo -- ocorre também na literatura de robótica e constitui tendência geral do campo. Mostra-se, por fim, que o sinal que cada estudo atribui ao efeito da tecnologia está associado ao tipo de evidência empregado, o que ajuda a explicar parte da aparente falta de convergência do campo.

% Em terceiro lugar, há a relevância do tema para o Brasil. As projeções e  os achados que dominam a literatura provêm majoritariamente de economias avançadas, mas seus efeitos sobre uma economia em desenvolvimento -- com estrutura ocupacional, informalidade e inserção em plataformas digitais próprias -- não são imediatamente transponíveis. Sistematizar o que se sabe é  condição necessária para discutir, com a devida cautela, o que esses achados podem significar no contexto brasileiro, tema retomado no Capítulo~\ref{cap:brasil}.


Além desta introdução, o trabalho organiza-se em cinco capítulos. O  Capítulo~\ref{cap:refteo} constrói o referencial teórico, percorrendo o caminho que vai da mudança técnica enviesada por qualificação à abordagem de tarefas e ao arcabouço de Acemoglu e Restrepo, cujos quatro mecanismos -- deslocamento, reinstalação, complementaridade e demanda agregada -- estruturam toda a análise.  O Capítulo~\ref{cap:metodologia} detalha o desenho da revisão sistemática: os critérios de inclusão, a estratégia de busca, o processo de seleção segundo o protocolo \gls{prisma}, o esquema de extração de dados e a avaliação de qualidade.

O Capítulo~\ref{cap:resultados} apresenta os resultados em três seções. A Seção~\ref{cap:descritivas} caracteriza o corpus; a Seção~\ref{cap:janelas} faz a leitura longitudinal por janela tecnológica; e a Seção~\ref{cap:comparacao}, núcleo do trabalho, conduz a comparação pareada entre os períodos anterior e posterior ao ChatGPT, distinguindo continuidades e rupturas. 
O Capítulo~\ref{cap:discussao} os confronta com o referencial teórico, examinando se o deslocamento do diagnóstico acompanha a data ou o objeto tecnológico dos estudos e quão sensíveis são as conclusões às escolhas metodológicas. Por fim, o Capítulo~\ref{cap:fim} retoma os principais achados, reconhece as limitações da revisão e aponta direções para pesquisa futura.
